% manuscript-snippet.tex
% Stripped-down working paper substrate for /review-paper --peer AER demo.
% Approx 3 pages of LaTeX. Sanitized — staggered DiD on Medicaid expansion.

\documentclass[11pt]{article}
\usepackage[margin=1in]{geometry}
\usepackage{natbib}

\title{Public Health Insurance and Small Business Formation: \\
       Evidence from Staggered Medicaid Expansion}
\author{[Author 1] \and [Author 2]}
\date{April 2026}

\begin{document}
\maketitle

\begin{abstract}
We examine how Medicaid expansion under the Affordable Care Act affected new
small-business formation across US counties. Using a staggered difference-in-differences
design that compares counties in expanding states to counties in never-expanding states,
we estimate that Medicaid expansion increased the small-business formation rate by
approximately 3.2 percentage points over a five-year horizon, with the strongest effects
concentrated in counties with above-median uninsurance rates pre-expansion. Our estimates
are robust to alternative comparison-group constructions, exclusion of border counties,
and the Callaway-Sant'Anna (2021) heterogeneity-robust ATT(g,t) estimator. We interpret
these results as evidence that public health insurance reduces a binding entrepreneurship
constraint, with implications for the design of safety-net programs.
\end{abstract}

\section{Introduction}

A long-standing question in labor economics asks whether the link between health
insurance and employment in the US distorts entrepreneurship choices \citep{MaduraEtAl2014,
GarthwaiteGrossNotowidigdo2014}. The theoretical mechanism — \emph{employment lock} —
predicts that workers stay in salaried jobs to retain employer-sponsored insurance,
suppressing transitions into self-employment.

In this paper, we exploit the staggered timing of state Medicaid expansions under the
2014 Affordable Care Act, which provided non-employer-tied health insurance to non-elderly
adults below 138\% of the federal poverty line. We examine whether expansion increased
the rate of new small-business formation in expanding counties relative to non-expanding
counties.

[Standard introduction discussion: relevance, contribution, related literature, roadmap.]

\section{Empirical Strategy}

\subsection{Sample Construction}

Our analysis sample consists of US counties from 2009 to 2022. We restrict the sample to
counties with at least five years of pre-period data and observable Census Business
Formation Statistics (BFS) coverage, yielding a final sample of \textbf{14,562 county-year}
observations across 2,891 counties.

\subsection{Identification}

We use a staggered difference-in-differences design comparing counties in expanding states
to counties in never-expanding states. The identifying assumption is that, absent
expansion, treated and control counties would have followed parallel trends in
small-business formation rates.

[Pre-trends figure would go here in the real paper. We discuss event-study evidence
in Section 4.1, but the figure itself is in the supplementary material.]

\subsection{Estimator}

Our preferred specification uses the Callaway-Sant'Anna (2021) heterogeneity-robust
ATT(g,t) estimator with never-treated counties as the comparison group:
\[ \text{ATT}(g,t) = E[Y_t(g) - Y_t(\infty) \mid G = g] \]
where $G = g$ denotes treatment cohort year and $Y_t(\infty)$ is the never-treated
counterfactual. We aggregate ATT(g,t) into an event-study at horizon $\tau = t - g$
and report the average over post-treatment horizons.

For comparison, we also report two-way fixed-effects (TWFE) estimates with state and
year fixed effects, recognizing the now-well-known limitations under heterogeneous
treatment effects \citep{deChaisemartin2020,Goodmanbacon2021}.

\section{Main Results}

[Table 2 from the paper would appear here. Headline: a \textbf{3.2 percentage-point}
increase in small-business formation rate over a five-year horizon
(SE 0.083, $N$ = 14,562, $R^2$ = 0.31). The effect is statistically significant at
conventional levels and robust across the alternative comparison groups discussed
in Section 4.2.]

[Approx 1.5 pages of robustness, mechanism discussion, heterogeneity by pre-expansion
uninsurance rate. Discussion of why the effect is largest in high-uninsurance counties.]

\section{Conclusion}

We provide quasi-experimental evidence that Medicaid expansion increased small-business
formation rates in expanding counties, particularly those with high pre-expansion
uninsurance. The implications extend to broader debates over the design of social
safety nets and their interaction with entrepreneurship. Our results suggest that the
US system's tight coupling of health insurance and employment imposes real distortions
on the small-business sector.

% \bibliographystyle{aer}
% \bibliography{Bibliography}
% (Bibliography not shipped — references inline in the snippet text.)

\end{document}
